The fibrewise assertions proved above for
$b\in\cB(\CC)$ extend to all geometric fibres without an implicit
genericity step.  The locus of $b\in\cB$ for which $X_b$ is geometrically
integral is open for this proper flat family
\cite[IV$_3$, Th\'eor\`eme~12.2.4\textup{(viii)}]{EGA4}; it contains every
closed point by \cref{prop:theta-resolution}, hence equals $\cB$ because
$\cB$ is Jacobson.  Likewise, the nonsmooth locus of
$(\cX\setminus0)\to\cB$ is closed and would contain a closed point if
nonempty, while smoothness along the zero section would be an open condition
and would again contain a closed point if it occurred anywhere.  Thus every
geometric fibre is integral with singular locus exactly the origin, and the
relative smooth locus is
\[
 \cM=(\cX/\cB)_{\mathrm{sm}}=\cX\setminus0.
\]
The morphism $\cM\to\cB$ is smooth and surjective with geometrically
integral fibres.  Since $\cB$ is integral, $\cM$ is integral.

Let $\cN$ be the relative norm fibre in $\cC^{(3)}$.  Every Abel class on
$\cM$ has $h^0=1$, and the relative Abel map identifies
$\cN\times_{\cX}\cM$ with $\cM$.  The universal divisor on the relative
symmetric cube therefore gives a universal effective divisor of degree three
over $\cM$.

\subsection{A dominant normalized chart}

\begin{lemma}[Dominant normalized chart]\label[lemma]{lem:dominant-chart}
There is a nonempty open subscheme
\[
 T_{\QQ}\subset
 \operatorname{Spec}\QQ[c,c^{-1},m,m^{-1},e_1,e_2,e_3],
\]
hence an integral five-dimensional rational scheme, such that with
$\lambda=m^{-3}$, after base change to $\CC$ suitable dense opens of
\[
 T=T_{\QQ}\times_{\QQ}\CC
 \qquad\text{and}\qquad
 \cM
\]
are isomorphic.  On these opens the universal divisor is described by
\eqref{eq:chart-curve-divisor} below.  On $e_1\ne0$, put
\begin{equation}
 \delta=\frac{c}{e_1^3},\qquad
 \Lambda=e_1^3\lambda,\qquad
 E_2=\frac{e_2}{e_1^2},\qquad
 E_3=\frac{e_3}{e_1^3},\qquad
 Z=e_1z.
\label{eq:normalized-variables}
\end{equation}
The residual-octic family and its factor-ordering cover are defined over
$T_{\QQ}$ and are birationally pulled back from the normalized chart
$e_1=1$.
\end{lemma}

\begin{proof}
Let $D$ be the unique reduced degree-three divisor representing a point of a
dense open in $\cM$.  Its norm on $E$ is a unique nonvertical line
\[
 y=mx+n,
 \qquad m\ne0.
\]
Set
\[
 x'=mx+n-c,
 \qquad
 Y'=y-c,
 \qquad
 w'=w.
\]
Then the norm line is $Y'=x'$.  If the three selected lifts have
$w'$-coordinates $\alpha,\beta,\gamma$, let $e_i$ be their elementary
symmetric functions and write
\[
 p(W)=W^3-e_1W^2+e_2W-e_3.
\]
Newton's identities give
\[
 \begin{aligned}
 \alpha^3+\beta^3+\gamma^3
   &=e_1^3-3e_1e_2+3e_3=: \Sigma_1,\\
 (\alpha\beta)^3+(\alpha\gamma)^3+(\beta\gamma)^3
   &=e_2^3-3e_1e_2e_3+3e_3^2=: \Sigma_2.
 \end{aligned}
\]
Consequently
\[
 \prod_{\xi\in\{\alpha,\beta,\gamma\}}(x'-\xi^3)
 =(x')^3-\Sigma_1(x')^2+\Sigma_2x'-e_3^3.
\]
Since $\lambda=m^{-3}$, the curve and the selected divisor are reconstructed
in the primed coordinates by
\begin{equation}
 \begin{aligned}
 (Y'+c)^2
   &=G(x')\\
   &:=(x'+c)^2+\lambda\bigl((x')^3-\Sigma_1(x')^2
                  +\Sigma_2x'-e_3^3\bigr),\\
 D&:\quad x'=Y'=(w')^3,\qquad p(w')=0.
 \end{aligned}
\label{eq:chart-curve-divisor}
\end{equation}
Conversely these formulas recover the original short-Weierstrass model.
Indeed write
\[
 g_2=1-\lambda\Sigma_1,
 \qquad
 g_1=2c+\lambda\Sigma_2,
 \qquad
 g_0=c^2-\lambda e_3^3,
 \qquad
 n_0=-\frac{g_2}{3\lambda}.
\]
Since $m^3\lambda=1$, the substitution $x'=mx+n_0$ gives
\[
 (Y'+c)^2=x^3+Ax+B,
\]
where
\[
 A=m\left(g_1-\frac{g_2^2}{3\lambda}\right),
 \qquad
 B=G(n_0).
\]
The original norm line is then $y=mx+n_0+c$.  These constructions are
inverse on dense opens.  They are defined over $\QQ$, and therefore
construct $T_{\QQ}$ with
\begin{equation}
 \CC(T)=\CC(\cM).
\label{eq:chart-function-field}
\end{equation}

The equations are equivariant for
\[
 (x',Y',w',c)\longmapsto(r^3x',r^3Y',rw',r^3c),
\]
which in the original coordinates is the usual Weierstrass scaling
\[
 (x,y,w,c,A,B)
 \longmapsto
 (r^2x,r^3y,rw,r^3c,r^4A,r^6B).
\]
Thus
\[
 (c,\lambda,e_1,e_2,e_3,m)
 \longmapsto
 (r^3c,r^{-3}\lambda,re_1,r^2e_2,r^3e_3,rm).
\]
The basis of $H^0(C,3\kappa-D)$ used below has weights $3,3,4$, so its
projective coordinates transform by
\begin{equation}
 [u:v:z]\longmapsto[u:v:r^{-1}z].
\label{eq:section-scaling}
\end{equation}
Taking $r=e_1^{-1}$ gives \eqref{eq:normalized-variables} and
\[
 \QQ(c,\lambda,e_1,e_2,e_3)
 =\QQ(\delta,\Lambda,E_2,E_3)(e_1).
\]
If $q_0=m/e_1$, then
\begin{equation}
 \QQ(T_{\QQ})
 =\QQ(\delta,\Lambda,E_2,E_3)(q_0,e_1),
 \qquad q_0^3=\Lambda^{-1}.
\label{eq:cubic-chart-field}
\end{equation}
This is the required normalized description.
\end{proof}

\subsection{The norm conic}

All formulas in this subsection are first established over the rational
function fields indicated below.  Put $Y=y-c=w^3$.  Since
$\deg(3\kappa)=9>2g(C)-2$, Riemann--Roch gives
$h^0(C,3\kappa)=6$.  For the evaluation map
\[
 H^0(C,3\kappa)\longrightarrow H^0(D,(3\kappa)|_D),
\]
the kernel has the following basis:
\[
 s_0=x-Y,
 \qquad
 s_1=Y-e_1w^2+e_2w-e_3,
 \qquad
 s_2=(x-e_3)w-e_1Y+e_2w^2.
\]
Here is a quick check that no dimension count is hidden in this choice.  On
$D$ we have $x=Y=w^3$ and
$p(w)=w^3-e_1w^2+e_2w-e_3=0$, so all three sections vanish on $D$.  Moreover
$\deg(3\kappa-D)=6$ and $K_C=2\kappa$, hence Riemann--Roch gives
\[
 h^0(C,3\kappa-D)-h^0(C,D-\kappa)=3.
\]
On the chart under consideration
$[\mathcal O_C(D)\otimes\kappa^{-1}]\in X_{\mathrm{reg}}=X\setminus\{0\}$,
so the degree-zero line bundle $\mathcal O_C(D-\kappa)$ is nontrivial and
has no section.  Thus $h^0(C,3\kappa-D)=3$.  Finally, in the expansion over
$k(E)$ with basis $1,w,w^2$, the coefficient of $x$ in the $w$-part first
forces the coefficient of $s_2$ in a constant linear relation to vanish;
the coefficient of $x$ in the $1$-part then forces that of $s_0$ to vanish,
and the coefficient of $s_1$ follows.  The three displayed sections are
therefore a basis.

For $s=us_0+vs_1+zs_2$, write
\[
 s=\alpha+\beta w+\gamma w^2,
\]
where
\[
 \begin{aligned}
 \alpha&=u(x-Y)+v(Y-e_3)-e_1zY,\\
 \beta&=e_2v+z(x-e_3),\\
 \gamma&=-e_1v+e_2z.
 \end{aligned}
\]
Since $w^3=Y$, the cyclic norm is
\begin{equation}
 \Nm(s)=\alpha^3+\beta^3Y+\gamma^3Y^2-3\alpha\beta\gamma Y.
\label{eq:cyclic-norm}
\end{equation}
It is divisible on $E$ by the fixed norm line $Y-x$.  The quotient is a
section of $\mathcal O_E(6O)$ and hence has a unique expression
\[
 Q_s=q_{00}+q_{01}x+q_{02}Y+q_{11}x^2+q_{12}xY+q_{22}Y^2.
\]
Let
\[
 M_s=
 \begin{pmatrix}
 q_{00}&q_{01}/2&q_{02}/2\\
 q_{01}/2&q_{11}&q_{12}/2\\
 q_{02}/2&q_{12}/2&q_{22}
 \end{pmatrix}.
\]

\begin{proposition}[Determinant factorization]
\label[proposition]{prop:determinant-factorization}
Put
\[
 K_0=\QQ(c,\lambda,e_1,e_2,e_3).
\]
In $K_0[u,v,z]$ one has
\begin{equation}
 \det(M_s)=(u-v)F_8,
\label{eq:determinant-factorization}
\end{equation}
where $F_8$ is homogeneous of degree eight and $u-v\nmid F_8$.
\end{proposition}

